\section{Analysis Samples}


The studies reported here includes data logged with 
the FHC $\nu$ beam runs (%{\color{blue}
October 2012 to February 2013%}
)
and the  RHC $\overline{\nu}$ beam runs
(May 2014 to June 2014).

\subsection{Data samples and detector configuration}
The total PoT exposure where all detector data quality checks were passed 
for the FHC runs was 16.24$\times10^{19}$ 
and the corresponding total PoT exposure for the RHC runs was 4.30$\times 10^{19}$.
These integrated rates corresponds to roughly
$0.28\times10^{12}$ neutrinos and $0.06\times10^{12}$ antineutrinos
per cm$^{2}$ per 50 MeV at 0.6 GeV. 
The data samples in this paper used the available neutrino and antineutrino beam 
data taken when the P$\emptyset$D target bags were filled with water.  

\subsection{Monte Carlo simulation}

The analysis used simulated MC
samples with different beam and detector configurations for each data taking period. 
The simulations include the following: %two (three) main parts: 
\begin{enumerate}
\item
Secondary pions and kaons are produced in the graphite target and propagated through 
the magnetic horns into a helium filled pipe where 
they decay.  Secondary neutrinos and antineutrinos are created
and their fluxes and energy spectra 
are extrapolated to the near and far detectors.
%of the particles leaving the target complex 
\item
The neutrino and antineutrino interactions in the ND280 sub-detectors
were determined by the NEUT \citep{NEUT-1} MC generator 
that was used to calculate the interaction cross sections and 
the final state particle kinematics.
% for the different flavor neutrinos.
% in the ND280 detector. 
%For this part the NEUT MC generator is utilized.
%propagated through the ND280 detector complex.
\item
The detector simulation uses GEANT to propagate the final state particles 
%traving in ND280 sub detectors.
through the ND280 sub-detectors.
\end{enumerate}

%\subsubsection{Beam Flux Prediction}
%
%\subsubsection{(Anti-) Neutrino Interaction Model}
%Given a neutrino energy and a detector geometry, NEUT
%event generator determines the interaction mode of an event and calculates
%the kinematics of the final-state particles. It also simulates
%FSIs for hadrons as they traverse the target nucleus. The
%following interaction modes are provided for both CC and
%neutral-current (NC) interactions by NEUT:
%\begin{enumerate}
%\item
%Quasielastic scattering (CCQE or NCQE),
%\item
%Resonant pion production (CC1$\pi$ or NC1$\pi$),
%\item
%Deep inelastic scattering (CCDIS or NCDIS), and
%\item
%Coherent pion production.
%\end{enumerate}
%%Here N and N' denote nucleons, l is the lepton, and A is
%%the target nucleus.
%
%The utilized NEUT generator also included: 
%Spectral function for CCQE (A. Ankowski's paper), 
%Multi-nucleon correlations based on Nieves model (meson exchange current) 
%and Improved form factor in the Rein-Sehgal model for 1 $\pi$ production 
%(K.M. Graczyk, J.T. Sobczyk and J.A. Nowak's paper).
%
%From tuning to Super-Kamiokande atmospheric data and
%K2K data, the nominal axial mass $M^{QE}_A$ was set to 1.21 GeV. 
%
%Neutrino-induced pion production is simulated
%based on the Rein Sehgal model [57] 
%{
%[57] D. Rein and L. M. Sehgal, Ann. Phys. (N.Y.) 133, 79 (1981).
%} 
%in NEUT with the
%axial mass $M^{RES}_A$ = 1.21 GeV. 
%The parton distribution
%function GRV98 [58] 
%{
%[58] M. Gluck, E. Reya, and A. Vogt, Eur. Phys. J. C 5, 461 (1998).
%}
%with corrections by Bodek and Yang [59]
%{
%[59] A. Bodek and U. K. Yang, AIP Conf. Proc. 670, 110 (2003).
%}
%is used for the deep inelastic scattering interactions.
%Secondary interactions of pions inside the nucleus
%[so-called final state interactions (FSIs)] are simulated
%using an intranuclear cascade model based on the method
%of Oset [60]
%{
%[60] L. L. Salcedo, E. Oset, M. J. Vicente-Vacas, and C. Garcia-
%Recio, Nucl. Phys. A484, 557 (1988).
%}
%, tuned to external $\pi-^{12}C$ data.
